\documentclass[10pt,a4paper]{amsart}
\usepackage{amsmath,amssymb,amsthm,graphicx}
\usepackage{verbatim}
\newenvironment{reference}{\comment}{\endcomment}
%\newenvironment{reference}{}{}
\newenvironment{slogan}{\comment}{\endcomment}
\newenvironment{history}{\comment}{\endcomment}
% For commutative diagrams we use Xy-pic
\usepackage[all]{xy}

% We use 2cell for 2-commutative diagrams.
\xyoption{2cell}
% Theorem environments.
%
\theoremstyle{plain}
\newtheorem{theorem}[subsection]{Theorem}
\newtheorem{proposition}[subsection]{Proposition}
\newtheorem{lemma}[subsection]{Lemma}

\theoremstyle{definition}
\newtheorem{definition}[subsection]{Definition}
\newtheorem{example}[subsection]{Example}
\newtheorem{exercise}[subsection]{Exercise}
\newtheorem{situation}[subsection]{Situation}

\theoremstyle{remark}
\newtheorem{remark}[subsection]{Remark}
\newtheorem{remarks}[subsection]{Remarks}

\numberwithin{equation}{subsection}

% Macros
%
\def\lim{\mathop{\mathrm{lim}}\nolimits}
\def\colim{\mathop{\mathrm{colim}}\nolimits}
\def\Spec{\mathop{\mathrm{Spec}}}
\def\Hom{\mathop{\mathrm{Hom}}\nolimits}
\def\Ext{\mathop{\mathrm{Ext}}\nolimits}
\def\SheafHom{\mathop{\mathcal{H}\!\mathit{om}}\nolimits}
\def\SheafExt{\mathop{\mathcal{E}\!\mathit{xt}}\nolimits}
\def\Sch{\mathit{Sch}}
\def\Mor{\mathop{\mathrm{Mor}}\nolimits}
\def\Ob{\mathop{\mathrm{Ob}}\nolimits}
\def\Sh{\mathop{\mathit{Sh}}\nolimits}
\def\NL{\mathop{N\!L}\nolimits}
\def\CH{\mathop{\mathrm{CH}}\nolimits}
\def\proetale{{pro\text{-}\acute{e}tale}}
\def\etale{{\acute{e}tale}}
\def\QCoh{\mathit{QCoh}}
\def\Ker{\mathop{\mathrm{Ker}}}
\def\Im{\mathop{\mathrm{Im}}}
\def\Coker{\mathop{\mathrm{Coker}}}
\def\Coim{\mathop{\mathrm{Coim}}}

% Boxtimes
%
\DeclareMathSymbol{\boxtimes}{\mathbin}{AMSa}{"02}

%
% Macros for moduli stacks/spaces
%
\def\QCohstack{\mathcal{QC}\!\mathit{oh}}
\def\Cohstack{\mathcal{C}\!\mathit{oh}}
\def\Spacesstack{\mathcal{S}\!\mathit{paces}}
\def\Quotfunctor{\mathrm{Quot}}
\def\Hilbfunctor{\mathrm{Hilb}}
\def\Curvesstack{\mathcal{C}\!\mathit{urves}}
\def\Polarizedstack{\mathcal{P}\!\mathit{olarized}}
\def\Complexesstack{\mathcal{C}\!\mathit{omplexes}}
% \Pic is the operator that assigns to X its picard group, usage \Pic(X)
% \Picardstack_{X/B} denotes the Picard stack of X over B
% \Picardfunctor_{X/B} denotes the Picard functor of X over B
\def\Pic{\mathop{\mathrm{Pic}}\nolimits}
\def\Picardstack{\mathcal{P}\!\mathit{ic}}
\def\Picardfunctor{\mathrm{Pic}}
\def\Deformationcategory{\mathcal{D}\!\mathit{ef}}
\usepackage{cleveref}
\crefname{originallemma}{original lemma}{original lemmas}
\Crefname{originallemma}{Original lemma}{Original lemmas}
\makeatletter
\newlabel{lemma-base-change-quasi-compact}{{lemma-base-change-quasi-compact}{source}{}{sourcepointer-1}{}}
\newlabel{lemma-base-change-quasi-compact@cref}{{[originallemma][0][]lemma-base-change-quasi-compact}{[1][0][]source}{}{}{}}
\makeatother
\crefname{originallemma}{original lemma}{original lemmas}
\Crefname{originallemma}{Original lemma}{Original lemmas}
\makeatletter
\newlabel{lemma-gerbe-bijection-points}{{lemma-gerbe-bijection-points}{source}{}{sourcepointer-2}{}}
\newlabel{lemma-gerbe-bijection-points@cref}{{[originallemma][0][]lemma-gerbe-bijection-points}{[1][0][]source}{}{}{}}
\makeatother
\crefname{originallemma}{original lemma}{original lemmas}
\Crefname{originallemma}{Original lemma}{Original lemmas}
\makeatletter
\newlabel{lemma-gerbe-diagonal-quasi-compact}{{lemma-gerbe-diagonal-quasi-compact}{source}{}{sourcepointer-3}{}}
\newlabel{lemma-gerbe-diagonal-quasi-compact@cref}{{[originallemma][0][]lemma-gerbe-diagonal-quasi-compact}{[1][0][]source}{}{}{}}
\makeatother
\crefname{originallemma}{original lemma}{original lemmas}
\Crefname{originallemma}{Original lemma}{Original lemmas}
\makeatletter
\newlabel{lemma-quasi-compact-permanence}{{lemma-quasi-compact-permanence}{source}{}{sourcepointer-4}{}}
\newlabel{lemma-quasi-compact-permanence@cref}{{[originallemma][0][]lemma-quasi-compact-permanence}{[1][0][]source}{}{}{}}
\makeatother
\crefname{originalproposition}{original proposition}{original propositions}
\Crefname{originalproposition}{Original proposition}{Original propositions}
\makeatletter
\newlabel{proposition-open-stratum}{{proposition-open-stratum}{source}{}{sourcepointer-5}{}}
\newlabel{proposition-open-stratum@cref}{{[originalproposition][0][]proposition-open-stratum}{[1][0][]source}{}{}{}}
\makeatother
\crefname{originalreference}{original reference}{original references}
\Crefname{originalreference}{Original reference}{Original references}
\makeatletter
\newlabel{spaces-properties-lemma-quasi-compact-quasi-separated-spectral}{{spaces-properties-lemma-quasi-compact-quasi-separated-spectral}{source}{}{sourcepointer-6}{}}
\newlabel{spaces-properties-lemma-quasi-compact-quasi-separated-spectral@cref}{{[originalreference][0][]spaces-properties-lemma-quasi-compact-quasi-separated-spectral}{[1][0][]source}{}{}{}}
\makeatother
\crefname{originalreference}{original reference}{original references}
\Crefname{originalreference}{Original reference}{Original references}
\makeatletter
\newlabel{stacks-properties-definition-reduced-induced-stack}{{stacks-properties-definition-reduced-induced-stack}{source}{}{sourcepointer-7}{}}
\newlabel{stacks-properties-definition-reduced-induced-stack@cref}{{[originalreference][0][]stacks-properties-definition-reduced-induced-stack}{[1][0][]source}{}{}{}}
\makeatother
\crefname{originalreference}{original reference}{original references}
\Crefname{originalreference}{Original reference}{Original references}
\makeatletter
\newlabel{stacks-properties-lemma-points-cartesian}{{stacks-properties-lemma-points-cartesian}{source}{}{sourcepointer-8}{}}
\newlabel{stacks-properties-lemma-points-cartesian@cref}{{[originalreference][0][]stacks-properties-lemma-points-cartesian}{[1][0][]source}{}{}{}}
\makeatother
\crefname{originalreference}{original reference}{original references}
\Crefname{originalreference}{Original reference}{Original references}
\makeatletter
\newlabel{stacks-properties-lemma-quasi-compact-stack}{{stacks-properties-lemma-quasi-compact-stack}{source}{}{sourcepointer-9}{}}
\newlabel{stacks-properties-lemma-quasi-compact-stack@cref}{{[originalreference][0][]stacks-properties-lemma-quasi-compact-stack}{[1][0][]source}{}{}{}}
\makeatother
\crefname{originalreference}{original reference}{original references}
\Crefname{originalreference}{Original reference}{Original references}
\makeatletter
\newlabel{stacks-properties-lemma-topology-points}{{stacks-properties-lemma-topology-points}{source}{}{sourcepointer-10}{}}
\newlabel{stacks-properties-lemma-topology-points@cref}{{[originalreference][0][]stacks-properties-lemma-topology-points}{[1][0][]source}{}{}{}}
\makeatother
\crefname{originalreference}{original reference}{original references}
\Crefname{originalreference}{Original reference}{Original references}
\makeatletter
\newlabel{topology-definition-spectral-space}{{topology-definition-spectral-space}{source}{}{sourcepointer-11}{}}
\newlabel{topology-definition-spectral-space@cref}{{[originalreference][0][]topology-definition-spectral-space}{[1][0][]source}{}{}{}}
\makeatother
\title{The topological space of a quasi-compact stack with quasi-compact diagonal is spectral}
\author{The Stacks Project Authors}
\date{stacks-project:stacks-morphisms.tex:local-snapshot-20261012; GNU FDL 1.2 or later}
\begin{document}
\maketitle
\noindent\textit{Editorial note.} All delivered mathematical text is the author's own, copied from the pinned Stacks Project chapter source sources/stacks-algebraic-stacks-morphisms/pinned-original-source.tex; no statement or proof step was rewritten, reordered or removed.\par\medskip
\noindent\textit{Editorial note.} Preamble provenance: the Stacks Project keeps one single shared preamble for every chapter; this card takes it directly from the staged sources/stacks-algebraic-stacks-morphisms/preamble.tex as three verbatim spans (lines 8-12, 14-18 and 154-217), so no derived file was written and nothing was added or redefined. The wrapper class and the base packages are supplied by the assembler.\par\medskip
\noindent\textit{Editorial note.} Bibliography provenance: neither the counted statement nor its proof contains a citation (this lane rescans the whole span before writing the card), so the card carries no bibliography and no entry was reconstructed.\par\medskip
\noindent\textit{Editorial note.} Selection provenance: this card was opened under the file-based control contract (CONTROL.md), with the mechanical pre-card gate run on its statement and proof interval before the card was written.\par\medskip
\noindent\textit{Allowed standard prerequisites (editorial).}\begin{itemize}
\item The presentation of a quasi-compact algebraic stack by an affine scheme, and the openness and surjectivity of the induced map on topological spaces. Invoked at the start of the counted proof (source lines 6941-6945) as pointers to results of the chapter on properties of stacks. Produces the presentation on which the basis of quasi-compact opens and the stability under intersections are checked.
\item The identification of fibre products of presentations with the intersection of the corresponding opens of the topological space. Invoked inside the counted proof (source lines 6955-6957). Shows that the intersection of two quasi-compact opens is again quasi-compact, using the quasi-compactness of the diagonal.
\item The existence of a dense open gerbe inside a reduced stack, the bijection of points for a gerbe, and the spectrality of the space of a quasi-compact quasi-separated algebraic space. Invoked in the second half of the counted proof (source lines 6989-7003) as pointers to results of this and other chapters. Reduces sobriety to the case of a gerbe over an algebraic space, where the space is already known to be spectral.
\end{itemize}\medskip
\ifdefined\hypertarget\hypertarget{sourcepointer-1}{}\fi
\noindent\textit{Source reference.} Original lemma lemma-base-change-quasi-compact: sources/stacks-algebraic-stacks-morphisms/pinned-original-source.tex, source lines 2062--2062. Navigation pointer only: the counted text refers to the Stacks item lemma-base-change-quasi-compact, which belongs to another chapter of the project and is not part of this card. The pointer carries the target's own label as its number because this chapter's staged snapshot does not contain the other chapter's tag table; its content is not reproduced here.\par\medskip
\ifdefined\hypertarget\hypertarget{sourcepointer-2}{}\fi
\noindent\textit{Source reference.} Original lemma lemma-gerbe-bijection-points: sources/stacks-algebraic-stacks-morphisms/pinned-original-source.tex, source lines 6699--6699. Navigation pointer only: the counted text refers to the Stacks item lemma-gerbe-bijection-points, which belongs to another chapter of the project and is not part of this card. The pointer carries the target's own label as its number because this chapter's staged snapshot does not contain the other chapter's tag table; its content is not reproduced here.\par\medskip
\ifdefined\hypertarget\hypertarget{sourcepointer-3}{}\fi
\noindent\textit{Source reference.} Original lemma lemma-gerbe-diagonal-quasi-compact: sources/stacks-algebraic-stacks-morphisms/pinned-original-source.tex, source lines 6998--6998. Navigation pointer only: the counted text refers to the Stacks item lemma-gerbe-diagonal-quasi-compact, which belongs to another chapter of the project and is not part of this card. The pointer carries the target's own label as its number because this chapter's staged snapshot does not contain the other chapter's tag table; its content is not reproduced here.\par\medskip
\ifdefined\hypertarget\hypertarget{sourcepointer-4}{}\fi
\noindent\textit{Source reference.} Original lemma lemma-quasi-compact-permanence: sources/stacks-algebraic-stacks-morphisms/pinned-original-source.tex, source lines 2082--2082. Navigation pointer only: the counted text refers to the Stacks item lemma-quasi-compact-permanence, which belongs to another chapter of the project and is not part of this card. The pointer carries the target's own label as its number because this chapter's staged snapshot does not contain the other chapter's tag table; its content is not reproduced here.\par\medskip
\ifdefined\hypertarget\hypertarget{sourcepointer-5}{}\fi
\noindent\textit{Source reference.} Original proposition proposition-open-stratum: sources/stacks-algebraic-stacks-morphisms/pinned-original-source.tex, source lines 6826--6826. Navigation pointer only: the counted text refers to the Stacks item proposition-open-stratum, which belongs to another chapter of the project and is not part of this card. The pointer carries the target's own label as its number because this chapter's staged snapshot does not contain the other chapter's tag table; its content is not reproduced here.\par\medskip
\ifdefined\hypertarget\hypertarget{sourcepointer-6}{}\fi
\noindent\textit{Source reference.} Original reference spaces-properties-lemma-quasi-compact-quasi-separated-spectral: sources/stacks-algebraic-stacks-morphisms/pinned-original-source.tex, source lines 7003--7003. Navigation pointer only: the counted text refers to the Stacks item spaces-properties-lemma-quasi-compact-quasi-separated-spectral, which belongs to another chapter of the project and is not part of this card. The pointer carries the target's own label as its number because this chapter's staged snapshot does not contain the other chapter's tag table; its content is not reproduced here.\par\medskip
\ifdefined\hypertarget\hypertarget{sourcepointer-7}{}\fi
\noindent\textit{Source reference.} Original reference stacks-properties-definition-reduced-induced-stack: sources/stacks-algebraic-stacks-morphisms/pinned-original-source.tex, source lines 6968--6968. Navigation pointer only: the counted text refers to the Stacks item stacks-properties-definition-reduced-induced-stack, which belongs to another chapter of the project and is not part of this card. The pointer carries the target's own label as its number because this chapter's staged snapshot does not contain the other chapter's tag table; its content is not reproduced here.\par\medskip
\ifdefined\hypertarget\hypertarget{sourcepointer-8}{}\fi
\noindent\textit{Source reference.} Original reference stacks-properties-lemma-points-cartesian: sources/stacks-algebraic-stacks-morphisms/pinned-original-source.tex, source lines 2126--2126. Navigation pointer only: the counted text refers to the Stacks item stacks-properties-lemma-points-cartesian, which belongs to another chapter of the project and is not part of this card. The pointer carries the target's own label as its number because this chapter's staged snapshot does not contain the other chapter's tag table; its content is not reproduced here.\par\medskip
\ifdefined\hypertarget\hypertarget{sourcepointer-9}{}\fi
\noindent\textit{Source reference.} Original reference stacks-properties-lemma-quasi-compact-stack: sources/stacks-algebraic-stacks-morphisms/pinned-original-source.tex, source lines 1935--1935. Navigation pointer only: the counted text refers to the Stacks item stacks-properties-lemma-quasi-compact-stack, which belongs to another chapter of the project and is not part of this card. The pointer carries the target's own label as its number because this chapter's staged snapshot does not contain the other chapter's tag table; its content is not reproduced here.\par\medskip
\ifdefined\hypertarget\hypertarget{sourcepointer-10}{}\fi
\noindent\textit{Source reference.} Original reference stacks-properties-lemma-topology-points: sources/stacks-algebraic-stacks-morphisms/pinned-original-source.tex, source lines 2128--2128. Navigation pointer only: the counted text refers to the Stacks item stacks-properties-lemma-topology-points, which belongs to another chapter of the project and is not part of this card. The pointer carries the target's own label as its number because this chapter's staged snapshot does not contain the other chapter's tag table; its content is not reproduced here.\par\medskip
\ifdefined\hypertarget\hypertarget{sourcepointer-11}{}\fi
\noindent\textit{Source reference.} Original reference topology-definition-spectral-space: sources/stacks-algebraic-stacks-morphisms/pinned-original-source.tex, source lines 6960--6960. Navigation pointer only: the counted text refers to the Stacks item topology-definition-spectral-space, which belongs to another chapter of the project and is not part of this card. The pointer carries the target's own label as its number because this chapter's staged snapshot does not contain the other chapter's tag table; its content is not reproduced here.\par\medskip
\begin{lemma}
\label{lemma-spectral-qc-diagonal-qc}
Let $\mathcal{X}$ be a quasi-compact algebraic stack
whose diagonal $\Delta$ is quasi-compact.
Then $|\mathcal{X}|$ is a spectral topological space.
\end{lemma}

\begin{proof}
Choose an affine scheme $U$ and a surjective smooth morphism
$U \to \mathcal{X}$, see
Properties of Stacks, Lemma \ref{stacks-properties-lemma-quasi-compact-stack}.
Then $|U| \to |\mathcal{X}|$ is continuous, open, and surjective, see
Properties of Stacks, Lemma \ref{stacks-properties-lemma-topology-points}.
Hence the quasi-compact opens of $|\mathcal{X}|$ form a basis
for the topology. For $W_1, W_2 \subset |\mathcal{X}|$ quasi-compact open,
we may choose a quasi-compact opens $V_1, V_2$ of $U$ mapping
to $W_1$ and $W_2$. Since $\Delta$ is quasi-compact,
we see that
$$
V_1 \times_\mathcal{X} V_2 =
(V_1 \times V_2) \times_{\mathcal{X} \times \mathcal{X}, \Delta} \mathcal{X}
$$
is quasi-compact. Then image of $|V_1 \times_\mathcal{X} V_2|$ in
$|\mathcal{X}|$ is $W_1 \cap W_2$ by
Properties of Stacks, Lemma \ref{stacks-properties-lemma-points-cartesian}.
Thus $W_1 \cap W_2$ is quasi-compact.
To finish the proof, it suffices to show that $|\mathcal{X}|$
is sober, see Topology, Definition \ref{topology-definition-spectral-space}.

\medskip\noindent
Let $T \subset |\mathcal{X}|$ be an irreducible closed subset.
We have to show $T$ has a unique generic point.
Let $\mathcal{Z} \subset \mathcal{X}$ be the reduced induced
closed substack corresponding to $T$, see
Properties of Stacks, Definition
\ref{stacks-properties-definition-reduced-induced-stack}.
Since $\mathcal{Z} \to \mathcal{X}$ is a closed immersion,
we see that $\Delta_\mathcal{Z}$ is quasi-compact:
first show that $\mathcal{Z} \to \mathcal{X} \times \mathcal{X}$
is quasi-compact as the composition of $\mathcal{Z} \to \mathcal{X}$
with $\Delta$, then write $\mathcal{Z} \to \mathcal{X} \times \mathcal{X}$
as the composition of $\Delta_\mathcal{Z}$ and
$\mathcal{Z} \times \mathcal{Z} \to \mathcal{X} \times \mathcal{X}$ and
use Lemma \ref{lemma-quasi-compact-permanence}
and the fact that
$\mathcal{Z} \times \mathcal{Z} \to \mathcal{X} \times \mathcal{X}$
is separated. Thus we reduce to the case discussed in the next
paragraph.

\medskip\noindent
Assume $\mathcal{X}$ is quasi-compact, $\Delta$ is quasi-compact,
$\mathcal{X}$ is reduced, and $|\mathcal{X}|$ irreducible.
We have to show $|\mathcal{X}|$ has a unique generic point.
Since $\mathcal{I}_\mathcal{X} \to \mathcal{X}$ is a base change of $\Delta$,
we see that $\mathcal{I}_\mathcal{X} \to \mathcal{X}$
is quasi-compact (Lemma \ref{lemma-base-change-quasi-compact}).
Thus there exists a dense open substack $\mathcal{U} \subset \mathcal{X}$
which is a gerbe by
Proposition \ref{proposition-open-stratum}.
In other words, $|\mathcal{U}| \subset |\mathcal{X}|$
is open dense. Thus we may assume that $\mathcal{X}$ is a gerbe.
Say $\mathcal{X} \to X$ turns $\mathcal{X}$ into a gerbe over the
algebraic space $X$. Then $|\mathcal{X}| \cong |X|$ by
Lemma \ref{lemma-gerbe-bijection-points}.
In particular, $X$ is quasi-compact.
By Lemma \ref{lemma-gerbe-diagonal-quasi-compact}
we see that $X$ has quasi-compact diagonal,
i.e., $X$ is a quasi-separated algebraic space.
Then $|X|$ is spectral by
Properties of Spaces, Lemma
\ref{spaces-properties-lemma-quasi-compact-quasi-separated-spectral}
which implies what we want is true.
\end{proof}

\end{document}
